%
\usepackage{amsmath}
\usepackage{amssymb}
\usepackage{amsfonts}
\usepackage{amsthm}
\usepackage{mathtools}
\usepackage{bm}
\usepackage{bbm}
\usepackage{dsfont}
\usepackage{cases}
\usepackage{nicefrac}
%
\usepackage{bbding}
\usepackage{pifont}

%
\usepackage[utf8]{inputenc} %
\usepackage[T1]{fontenc}    %
\usepackage{textcomp}

%
\usepackage{etoc}  %

%
\usepackage{graphicx}
\usepackage{wrapfig}
\usepackage{float}
\usepackage{caption}
\usepackage{subcaption}     %
\usepackage{array}
\usepackage{multirow}
\usepackage{multicol}
\usepackage{booktabs}
\usepackage{makecell}
\usepackage{colortbl}

%
\usepackage{xr-hyper}
\usepackage{hyperref}       %
\usepackage[capitalise]{cleveref}
\crefname{equation}{}{}
\crefname{figure}{Fig.}{Figs.}
\crefname{section}{Sec.}{Secs.}
\crefname{appendix}{App.}{Apps.}
\crefname{table}{Tab.}{Tabs.}
\usepackage{url}

\usepackage{tikz}
\usetikzlibrary{positioning, calc, matrix} 

\usepackage[makeroom]{cancel}

%
%
%
%
%
%
%

\hypersetup{
  final,
  colorlinks,
  linkcolor=ourred,
  citecolor=ourblue,
  urlcolor=url,
}

%
\usepackage{algorithm}
\usepackage{algpseudocode}
%
%

%
\usepackage[page,toc,titletoc,title]{appendix}
\usepackage{microtype}
\usepackage{xargs}
\usepackage{xspace}
\usepackage[normalem]{ulem}
\usepackage{soul}
\usepackage{fancyvrb}

%
\usepackage{lipsum}
\usepackage{blindtext}

%
\usepackage{xcolor}
\definecolor{ourblue}{rgb}{0.368,0.507,0.71}
\definecolor{ourorange}{rgb}{0.881,0.611,0.142}
\definecolor{ourgreen}{rgb}{0.56,0.692,0.195}
\definecolor{ourgreen2}{rgb}{0.46,0.792,0.195}
\definecolor{ourred}{rgb}{0.923,0.386,0.209}
\definecolor{ourviolet}{rgb}{0.528,0.471,0.701}
\definecolor{ourbrown}{rgb}{0.772,0.432,0.102}
\definecolor{ourlightblue}{rgb}{0.364,0.619,0.782}
\definecolor{ourbrightblue}{RGB}{93,135,237}
\definecolor{ourdarkgreen}{rgb}{0.572,0.586,0.}

\definecolor{ourcyan2}{rgb}{0.125,0.722,0.804}
\definecolor{ourred2}{rgb}{0.863,0.184,0.047}
\definecolor{ouryellow2}{cmyk}{0,0.16,1.0,0.07}
\definecolor{ourviolet2}{cmyk}{0.55,0.56,0,0.47}
\definecolor{ourorange2}{cmyk}{0,0.46,0.89,0.11}
\definecolor{grayseq}{RGB}{120,120,120}

\definecolor{discretecolor}{RGB}{11,83,150}
\definecolor{gaussiancolor}{RGB}{230,145,56}
\definecolor{argmaxcolor}{RGB}{154,0,0}

\definecolor{customred}{RGB}{169, 45, 59}
\definecolor{url}{HTML}{d95225}

\definecolor{olivegreen}{RGB}{165,185,115}
\definecolor{olivegreendark}{RGB}{145,165,95}

\definecolor{maroon}{RGB}{140, 45, 45}

\usepackage[most]{tcolorbox} %
\usepackage{soul}
\newcommand{\mathcolorbox}[2]{\colorbox{#1}{$\displaystyle #2$}}
\newcommand{\hlfancy}[2]{\sethlcolor{#1}\hl{#2}}

%
\makeatletter
\DeclareRobustCommand\onedot{\futurelet\@let@token\@onedot}
\def\@onedot{\ifx\@let@token.\else.\null\fi\xspace}
\def\eg{e.g\onedot}
\def\ie{i.e\onedot}
\def\cf{cf\onedot}
\def\etc{etc\onedot}
\def\wrt{w.r.t\onedot}
\makeatother

\newcommand{\Real}{\ensuremath{\mathbb R}}        %
\newcommand{\Unit}{\ensuremath{\mathbb I}}        %
\newcommand{\T}{\ensuremath{\top}}                %
\newcommand{\bfI}{\mathbf{I}}

%

%
%
%
%
%
%
%
%

\newtcolorbox{corollarybox}{
  breakable,
  enhanced,
  colback=ourbrightblue!12,
  colframe=ourbrightblue!80!black,
  left=1pt,
  right=1pt,
  top=1pt,
  bottom=1pt,
  boxrule=1pt
}

%
%
%
%
%
%
%
%

%
%
%
%
%
%
%
%
%
%
%

\newtcolorbox[auto counter]{findingbox}[1][]{
  breakable,
  enhanced,
  colback=ourbrightblue!8,
  colframe=ourbrightblue!80!black,
  left=2pt,
  right=2pt,
  top=1pt,
  bottom=1pt,
  boxrule=1pt,
  before skip=0.7em,
  after skip=0.7em,
  before upper={\textbf{Finding~\thetcbcounter: } },
  #1
}

\newtcolorbox{samplebox}[1][]{
  breakable,
  enhanced,
  colback=ourorange!6,
  colframe=ourorange!75!black,
  left=6pt, right=6pt, top=4pt, bottom=4pt,
  boxrule=0.6pt,
  arc=2.5pt,
  fonttitle=\bfseries\small,
  coltitle=black,
  colbacktitle=ourorange!22,
  attach boxed title to top left={yshift=-2.5mm, xshift=5mm},
  boxed title style={size=small, sharp corners=all, boxrule=0.3pt, arc=2pt, colframe=ourorange!75!black},
  fontupper=\small,
  title=Sample,
  #1
}

\newcommand{\zhihan}[1]{\textcolor{blue}{[zhihan: #1]}}
\newcommand{\subham}[1]{\textcolor{red}{[subham: #1]}}
\newcommand{\yongxin}[1]{\textcolor{cyan}{[yongxin: #1]}}
\newcommand{\wei}[1]{\textcolor{orange}{[wei: #1]}}
\newcommand{\av}[1]{\textcolor{teal}{[av: #1]}}
\definecolor{justingreen}{rgb}{0.0078, 0.4431, 0.2823}
\newcommand{\justin}[1]{\textcolor{justingreen}{[justin: #1]}}
\newcommand{\Morteza}[1]{\textcolor{purple}{{\bf Morteza:}  #1}}
\newcommand{\john}[1]{\textcolor{magenta}{{\bf John:}  #1}}

\def\method{Eso-LMs}
\def\methoda{Eso-LMs (A)}
\def\methodbsingular{Eso-LM}
\def\methodb{Eso-LMs}
\def\model{DiTA}

\usepackage{pifont}
\newcommand{\cmark}{\textcolor{green!70!black}{\ding{51}}}
\newcommand{\ccross}{\textcolor{red!70!black}{\ding{55}}}


%
%

%
%
%
%

\definecolor{topone}{RGB}{112, 149, 239}   %
\definecolor{toptwo}{RGB}{174, 195, 246}   %
\definecolor{topthree}{RGB}{223, 231, 251} %

\newcommand{\topone}[1]{\cellcolor{topone}#1}
\newcommand{\toptwo}[1]{\cellcolor{toptwo}#1}
\newcommand{\topthree}[1]{\cellcolor{topthree}#1}

% 5-step blue gradient (dark=best) for ranking the top-5 cells
\definecolor{rankone}{RGB}{99, 137, 236}    %
\definecolor{ranktwo}{RGB}{140, 168, 242}   %
\definecolor{rankthree}{RGB}{174, 195, 246}  %
\definecolor{rankfour}{RGB}{201, 216, 249}  %
\definecolor{rankfive}{RGB}{224, 233, 252}  %
\newcommand{\rankone}[1]{\cellcolor{rankone}#1}
\newcommand{\ranktwo}[1]{\cellcolor{ranktwo}#1}
\newcommand{\rankthree}[1]{\cellcolor{rankthree}#1}
\newcommand{\rankfour}[1]{\cellcolor{rankfour}#1}
\newcommand{\rankfive}[1]{\cellcolor{rankfive}#1}

%
\usepackage{arydshln}

\usepackage{amsthm}
\usepackage{thmtools}
\usepackage{thm-restate}

\newtheorem{theorem}{Theorem}
\crefname{theorem}{Thm.}{Thms.}
\newtheorem{lemma}{Lemma}
\crefname{lemma}{Lem.}{Lems.}
\newtheorem{assumption}{Assumption}
\crefname{assumption}{Assump.}{Assumps.}
\newtheorem{proposition}{Proposition}
\crefname{proposition}{Prop.}{Props.}
\newtheorem{corollary}{Corollary}
\crefname{corollary}{Cor.}{Cors.}
\newtheorem{definition}{Definition}
\crefname{definition}{Def.}{Defs.}
\newtheorem{remark}{Remark}
\crefname{remark}{Rmk.}{Rmks.}
\crefname{algocf}{Alg.}{Algs.}
\crefname{algorithm}{Alg.}{Algs.}

\newenvironment{sketchofproof}{%
  \renewcommand{\qedsymbol}{}%
  \begin{proof}[Sketch of proof]%
}{%
  \end{proof}%
}

\usepackage{titlesec}

%

%
\titlespacing*{\section}{0pt}{0.5\baselineskip}{0.2\baselineskip}
\titlespacing*{\subsection}{0pt}{0.3\baselineskip}{0.1\baselineskip}
\titlespacing*{\subsubsection}{0pt}{0.2\baselineskip}{0.1\baselineskip}
\titlespacing*{\paragraph}{0pt}{1.5ex plus 0.5ex minus 0.2ex}{1em}

\newcommand{\replaidscx}{20}
\newcommand{\replaidnscx}{27}

\newcommand{\replaidscppl}{22.1}
\newcommand{\replaidnscppl}{23.6}
\usepackage{minted}
% \usepackage{xcolor}
\usemintedstyle{friendly}
%
%
%
%

%
%
%
%

\usepackage{listings}

% 2. Create a custom syntax highlighter
% 1. Define the exact colors from your first screenshot
\definecolor{elfcomment}{HTML}{3B9C9C} % Teal
\definecolor{elfblue}{HTML}{1A237E}    % Deep Blue
\definecolor{elfmagenta}{HTML}{C2185B} % Magenta/Pink
\definecolor{elforange}{HTML}{E65100}  % Orange

% 2. Create a custom syntax highlighter
\lstdefinelanguage{ELF}{
    language=Python,
    basicstyle=\ttfamily\small\color{black},
    commentstyle=\color{elfcomment},
    stringstyle=\color{elforange},
    columns=fullflexible,                  % <-- ADD THIS LINE
    keepspaces=true,                       % <-- ADD THIS LINE
    showstringspaces=false,
    % Keywords (Blue)
    keywordstyle=[1]\color{elfblue},
    morekeywords=[1]{if, else, for, in, range, len},
    % Custom Functions (Magenta)
    keywordstyle=[2]\color{elfmagenta},
    morekeywords=[2]{encode, uniform, sample_t, randn_like, net, mse_loss, corrupt, unembed, ce_loss, randn, argmax, adaptive_bs, gamma_tilde_net, gamma_bounds, sqrt, sigmoid, prior_loss},
    % Operators (Blue)
    % Note: The (t = 1) rule protects your comment from turning blue!
    literate=
     *{=}{{\textcolor{elfblue}{=}}}{1}
      {+}{{\textcolor{elfblue}{+}}}{1}
      {-}{{\textcolor{elfblue}{-}}}{1}
      {<}{{\textcolor{elfblue}{<}}}{1}
      {*}{{\textcolor{elfblue}{*}}}{1}
      {/}{{\textcolor{elfblue}{/}}}{1}
      {(t = 1)}{{\textcolor{elfcomment}{(t = 1)}}}{7}, 
}